\documentclass[11pt,a4paper]{article}

\usepackage[utf8]{inputenc}
\usepackage[T1]{fontenc}
\usepackage{lmodern}
\usepackage{graphicx}
\usepackage{booktabs}
\usepackage{amsmath}
\usepackage{hyperref}
\usepackage{natbib}
\usepackage[margin=2.5cm]{geometry}
\usepackage{xcolor}
\usepackage{float}
\usepackage{multirow}
\usepackage{array}
\usepackage{tabularx}
\usepackage{CJKutf8}

\bibliographystyle{plainnat}

\title{Lexical and Rhythmic Fingerprints of Japanese LLM-Generated Text:\\A Cross-Model, Two-Layer Analysis}
\author{Ken Imoto\\Propel-lab\\Tokyo, Japan\\\texttt{imoto@propel-lab.co.jp}\\\url{https://kenimoto.dev}}
\date{July 2026}

\begin{document}
\begin{CJK}{UTF8}{min}

\maketitle

\begin{abstract}
Which feature family carries the human/AI signal in Japanese LLM-generated text: word choice or rhythm? Practitioner tools claim that the ``AI smell'' of Japanese text resides in rhythm (monotone sentence-length variation) rather than vocabulary, but this claim has not been tested against lexical features under a controlled design. We compare lexical and rhythmic fingerprints on a corpus of 350 texts from 7 LLMs (Claude 3 Haiku, Claude Sonnet 4, Claude Opus 4, GPT-3.5 Turbo, GPT-4o, GPT-OSS 20B, Llama 3.2 1B) against 700 human technical-blog articles. We first rebuild the excess-vocabulary statistics of our prior study on a document-occurrence basis with Laplace smoothing, frequency ratio $r$ and gap $\delta$, document-level $\chi^2$ tests with BH-FDR correction, and bootstrap confidence intervals, following the design of Kobak et al.\ (2025). This re-analysis is also a correction: 35\% of the 651 words we previously reported as significant do not survive at the document level, because token-level $\chi^2$ tests treat repeated occurrences within one document as independent evidence. We then measure 13 interpretable rhythm metrics, including burstiness of sentence-length sequences and a mora-based approximation of phonological rhythm. Three findings emerge. (1)~Rhythmic monotonization is real and shared: AI text shows lower burstiness (Cohen's $d=-0.96$), lower paragraph-structure variation ($d=-0.67$), and lower sentence-length variation ($d=-0.56$), with the same direction in all 7 models. (2)~Contrary to the practitioner claim, lexical features dominate human/AI discrimination: cross-validated AUC is 0.998 for lexical vs.\ 0.897 for rhythm features (paired difference $-0.101$, 95\% CI $[-0.105, -0.097]$), and the gap persists under document-length residualization and outlier-model exclusion. (3)~At the multivariate level, the two families form a two-layer fingerprint: lexical profiles identify \emph{which} model wrote a text (92\% accuracy in 7-way model identification vs.\ 51\% for rhythm), whereas rhythm features shift in a direction common to all models, with model-dependent magnitude. Rhythm is a shared accent of machine writing; vocabulary is the signature of the individual model.
\end{abstract}

\section{Introduction}

Detection of machine-generated text usually asks a binary question---was this written by a model?---and answers it with an opaque score. A complementary question is \emph{what} makes machine text machine-like. For English, one well-established answer is lexical: LLMs overuse a measurable set of words (\textit{delve}, \textit{intricate}, \textit{commendable}), and this excess vocabulary is strong enough to estimate LLM involvement at the population level~\citep{kobak2025excess, liang2024monitoring, juzek2025delve}. For Japanese, a competing answer has emerged from practice: the \texttt{natural-japanese} lint tool argues that the ``AI smell'' (AI臭) of Japanese text lives in \emph{rhythm} (monotone sentence lengths, uniform paragraph structure) rather than in vocabulary~\citep{mizoguchi2026natural}. The two answers have never been placed in the same experiment.

This paper does so. Using the corpus of our prior excess-vocabulary study~\citep{imoto2026excess} (350 texts from 7 LLMs and 700 human technical-blog articles), we quantify both feature families under a shared evaluation protocol and test two hypotheses:

\begin{itemize}
    \item \textbf{H1 (discriminative power)}: the human/AI discriminative power of lexical and rhythm features differs; the practitioner claim predicts rhythm $>$ lexical.
    \item \textbf{H2 (fingerprint layering)}: rhythm fingerprints are homogeneous across models while lexical fingerprints differentiate them, forming a two-layer structure.
\end{itemize}

Before the comparison, we rebuild our lexical statistics from the ground up. A methods audit of our prior study against \citet{kobak2025excess} found that its per-word $\chi^2$ tests used token counts as independent trials. Tokens within a document are strongly correlated (one article that repeats a term dozens of times contributes dozens of pseudo-independent observations), so the tests were affected by pseudoreplication, which inflates $\chi^2$ statistics and deflates $p$-values. We re-estimate all word statistics on a document-occurrence basis with Laplace smoothing, dual effect-size indices (frequency ratio $r$ and frequency gap $\delta$), BH-FDR correction, and document-level bootstrap confidence intervals. The result is a correction we state plainly: of the 651 words our prior study reported as significant, only 424 (65.1\%) survive the document-based re-analysis, and on the AI-excess side only 237 of 458 (51.7\%) survive.

This paper makes four contributions:
\begin{enumerate}
    \item A \textbf{document-occurrence re-analysis} of Japanese excess vocabulary with $r$/$\delta$ indices, Laplace smoothing, FDR control, and bootstrap CIs, including a quantified correction of our prior token-based results.
    \item The first \textbf{controlled comparison of lexical vs.\ rhythmic discriminative power} for Japanese human/AI classification, with leakage-safe within-fold feature selection, document-length residualization, and outlier sensitivity analyses.
    \item A set of \textbf{13 interpretable rhythm metrics} for Japanese, including burstiness of sentence-length sequences in characters and in approximated moras, measured across 7 LLMs.
    \item Evidence for a \textbf{two-layer fingerprint structure}: rhythm shifts in a direction shared by all models (layer 1), while lexical profiles and sentence-ending rhetoric are model-specific (layer 2).
\end{enumerate}

\section{Related Work}

\paragraph{Detection of LLM-Generated Text.}
Machine-generated text (MGT) detection is by now a well-developed research area; \citet{wu2025survey} organize the field into watermarking, statistical zero-shot methods, fine-tuned classifiers, and human-assisted approaches. The statistical line began with GLTR~\citep{gehrmann2019gltr}, which exposes per-token probability, rank, and predictive entropy under a reference language model and raised human detection accuracy from 54\% to 72\%. DetectGPT~\citep{mitchell2023detectgpt} introduced zero-shot detection based on the observation that machine-generated text tends to occupy negative-curvature regions of a model's log-probability function, requiring no training data or model-specific classifier. Binoculars~\citep{hans2024binoculars} contrasts the perplexity assigned by two closely related pretrained models and reports detecting over 90\% of ChatGPT-generated samples at a 0.01\% false positive rate, without using any data from the target model. These detectors answer \emph{whether} a text is machine-generated, but they return an opaque score: they do not identify which lexical choices or which rhythmic properties make a text machine-like. Interpretable alternatives built on classical stylometric features such as perplexity, burstiness, lexical diversity, and readability have begun to be examined, e.g.\ to probe how diffusion-based generators evade detectors tuned to autoregressive output~\citep{tarim2025detection}. Our study addresses this explanatory side of the problem: quantifying \emph{which} feature families carry the human/AI signal in Japanese.

\paragraph{Excess Vocabulary and Lexical Fingerprints.}
A parallel line of work characterizes LLM output at the word level. \citet{kobak2025excess} formalized the \emph{excess vocabulary} approach: rather than assuming a word list a priori, they detect words whose frequency rose abnormally after the introduction of LLMs, analyzing over 15 million PubMed abstracts (2010--2024) and estimating that at least 13.5\% of 2024 abstracts were processed with LLMs.
% Verified against the published Science Advances version (Sci. Adv. 11(27):eadt3813, 2025): abstract states "over 15 million biomedical abstracts" and "at least 13.5% of 2024 abstracts were processed with LLMs".
\citet{liang2024monitoring} estimated LLM involvement in AI-conference peer reviews at the population level from disproportionately used adjectives (e.g., \textit{commendable}, \textit{meticulous}, \textit{intricate}), concluding that 6.5--16.9\% of reviews were substantially modified by LLMs. \citet{juzek2025delve} identified 21 focal words, including \textit{delve}, that occur at 3--5$\times$ human frequency in LLM output, systematically ruled out model architecture, training data, and decoding algorithms as explanations, and pointed instead to human feedback; a follow-up study supports this attribution with human-participant experiments~\citep{juzek2025overuse}. All of these studies concern English. In prior work, we transferred the excess-vocabulary methodology to Japanese, identifying 651 statistically significant excess words across 7 LLMs against a human technical-blog baseline~\citep{imoto2026excess}, and separately catalogued 16 structural patterns characteristic of Japanese AI text~\citep{imoto2025slop}. The present study reuses the corpus of the first study directly and, as detailed in Section~\ref{sec:reanalysis}, corrects its statistical basis.

\paragraph{Cross-Lingual and Japanese Studies.}
Multilingual MGT resources largely bypass Japanese. MULTITuDE covers 74{,}081 texts generated by 8 LLMs in 11 languages~\citep{macko2023multitude}, and M4 spans multiple generators, domains, and languages~\citep{wang2024m4}; neither includes Japanese, and both document poor detector generalization to unseen languages and generators. The closest lexical study is \citet{juzek2026lexical}, which measures AI-associated lexical shifts in 34 languages, including Japanese, by comparing GPT-4.1-family outputs with human text and validating the resulting word sets against News Crawl data (2020--2021 vs.\ 2023--2024). However, results for Japanese are reported only at the aggregate level, no Japanese word list is published, and the analysis relies on a single model family, precluding cross-model comparison. On the detection side, \citet{iwata2025jsai} adapted Binoculars to Japanese with a focal-loss refinement, reaching accuracy and F-scores of 0.94 or higher for texts of 200 characters or more; their setup uses a single generator (GPT-3.5 Turbo) against the OSCAR corpus, evaluates detection scores only without characterizing features, and explicitly notes the scarcity of prior work on Japanese. Japanese quantitative stylometry itself offers a long tradition of statistical authorship analysis: bunsetsu-level part-of-speech patterns provide vocabulary-independent syntactic fingerprints~\citep{jin2013bunsetsu}, integrated classifiers combine morpheme, particle-distribution, and comma-placement features for author identification~\citep{jin2014integrated}, and stylistic statistics such as sentence length and part-of-speech ratios discriminate text genres~\citep{ishida2004academic}. This tradition has recently been extended to AI-generated text, showing that stylometric features separate LLM-generated Japanese from human writing even where human judges fail~\citep{zaitsu2023chatgpt, zaitsu2025stylometry}. These works establish morpheme- and syntax-level stylistic fingerprints for Japanese, but none contrasts lexical against rhythmic feature families across multiple LLMs.

\paragraph{Rhythm, Diversity, and Homogenization.}
A separate body of evidence indicates that LLM output is \emph{rhythmically} more uniform than human writing. \citet{munozortiz2024contrasting} compared human news text with six LLMs across three model families and found that human sentence-length distributions are more dispersed and human vocabulary more diverse; notably, differences among LLMs were smaller than the human--LLM gap, suggesting that rhythm-level uniformity is shared across models. \citet{sourati2025shrinking} document linguistic homogenization in both observational data and controlled experiments, affecting lexical, syntactic, and stylistic dimensions. \citet{padmakumar2024diversity} showed in a controlled co-writing experiment that a feedback-tuned model (InstructGPT) significantly reduces content diversity among writers while a base model (GPT-3) does not, and detection-side work finds that RLHF systematically alters the statistical properties and detectability of generated text~\citep{rlhf2025detectability}; together with the human-feedback attribution of lexical overuse~\citep{juzek2025overuse}, these results triangulate alignment tuning as a driver of both lexical and rhythmic convergence. A recent review synthesizes these homogenization effects at the lexical, syntactic, and content levels~\citep{homogenizing2026review}. Practitioner tools operationalize the rhythm intuition directly. GPTZero popularized \emph{burstiness} (the within-document variation of perplexity) as a detection signal, on the premise that human writers vary sentence structure and vocabulary more than LLMs do~\citep{gptzero2023burstiness}. For Japanese, the \texttt{natural-japanese} lint~\citep{mizoguchi2026natural} measures sentence-length rhythm in moras, burstiness, the coefficient of variation of paragraph structure, and the repetition rate of the contrastive construction ではなく (\textit{de wa naku}) over a corpus of 137 human and 406 AI texts (7 models), and argues that the ``AI smell'' of Japanese text resides in rhythm rather than vocabulary. This claim rests on practical observation; it has not been tested in a controlled comparison against lexical features.

\paragraph{Positioning of This Work.}
The intersection of these strands is empty: no existing study contrasts lexical and rhythmic fingerprints of AI-generated Japanese under a controlled multi-LLM design. Lexical-fingerprint research is English-centric~\citep{kobak2025excess, juzek2025delve, liang2024monitoring}; the one cross-lingual study covering Japanese publishes no Japanese word list and uses a single model family~\citep{juzek2026lexical}; Japanese detection work is score-based and single-generator~\citep{iwata2025jsai}; and rhythm-level homogeneity has been documented for English~\citep{munozortiz2024contrasting, sourati2025shrinking} or asserted by practitioner tools without academic validation~\citep{mizoguchi2026natural}. Building on the corpus and the lexical catalog of our prior studies~\citep{imoto2026excess, imoto2025slop}, we compare the discriminative power of lexical and rhythm feature families for the human/AI distinction in Japanese (H1), and test whether the two families form a two-layer fingerprint structure in which rhythm features are homogeneous across models while lexical features differentiate them (H2).

\section{Data and Methods}

\subsection{Corpus}

We reuse the corpus of our prior study~\citep{imoto2026excess} unchanged (read-only). The AI side comprises 350 documents: 7 LLMs---Claude 3 Haiku (\texttt{claude-3-haiku-20240307}), Claude Sonnet 4 (\texttt{claude-sonnet-4-20250514}), Claude Opus 4 (\texttt{claude-opus-4-20250514}), GPT-3.5 Turbo, GPT-4o, GPT-OSS 20B (Swallow), and Llama 3.2 1B---each generating 10 Japanese technical-blog themes 5 times under a unified prompt at default temperature. The human side comprises 700 pre-LLM articles (2020--2022; Qiita 567, Zenn 133). The lexical analyses use all 350 AI and 700 human documents. The rhythm and discrimination analyses exclude 4 human documents with fewer than 5 sentences, leaving 1{,}046 documents (AI 350; human 696 = Qiita 565 + Zenn 131).

All analyses fix the random seed (20260717) and record dependency versions and the source-corpus commit hash in machine-readable metadata files alongside the results (Python 3.10.12, numpy 2.2.6, scipy 1.15.3, scikit-learn 1.7.2, pandas 2.3.3, mecab-python3 1.0.12, unidic-lite 1.0.8).

\subsection{Document-Occurrence Excess Vocabulary}
\label{sec:lexmethod}

Our prior study computed word frequencies over tokens (occurrences per total tokens) and tested each word with a $\chi^2$ test whose contingency table used token totals as trial counts. A methods audit against \citet{kobak2025excess} identified this as pseudoreplication: tokens within one document are strongly correlated, so treating each token as an independent trial inflates the effective sample size (AI 121{,}243 / human 573{,}279 tokens) and makes near-arbitrary frequency differences ``significant.'' \citet{kobak2025excess} avoid exactly this by counting \emph{document occurrences}: a word occurring in a document counts once, however often it repeats (a binary document $\times$ word matrix).

We therefore rebuild the statistics as follows. Tokenization is identical to the prior study (MeCab~\citep{kudo2006mecab} with UniDic-lite~\citep{unidic2021}, lemma forms); we verified compatibility by reproducing the prior token totals exactly. For each word $w$, let $a$ be the number of documents containing $w$ and $b$ the corpus document count. The smoothed document frequency is $p = (a+1)/(b+1)$ (Laplace smoothing, following Kobak et al.'s Methods), computed separately for the AI corpus ($p$, $b=350$) and the human corpus ($q$, $b=700$). We report both of Kobak et al.'s complementary effect-size indices:
\begin{equation}
    r = \frac{p}{q}, \qquad \delta = p - q,
\end{equation}
where $r$ (frequency ratio) is sensitive to excess in low-frequency words and $\delta$ (frequency gap) to excess in high-frequency words. Significance is assessed per word with a $\chi^2$ test on the document-count contingency table $[[a_{\text{AI}}, 350-a_{\text{AI}}], [a_{\text{H}}, 700-a_{\text{H}}]]$ (Yates continuity correction; Fisher's exact test when a marginal is zero), corrected with Benjamini--Hochberg FDR at $\alpha=0.05$~\citep{benjamini1995fdr} as the primary criterion and Bonferroni as a stricter secondary criterion. For each word we compute percentile 95\% confidence intervals for $r$ and $\delta$ by bootstrap over documents (both corpora resampled, $B=1{,}000$); these CIs are per-word and not adjusted for multiplicity.

Two deliberate departures from \citet{kobak2025excess} deserve emphasis. First, their baseline is a \emph{counterfactual time-series extrapolation} within a single evolving corpus, with conservative projection ($q \geq p_{-2}$); ours is a cross-sectional comparison against a static pre-LLM human corpus, which is the appropriate design for controlled generation but has weaker confound separation (Section~\ref{sec:limitations}). Second, Kobak et al.\ perform \emph{no per-word hypothesis tests}: their robustness comes from effect-size thresholds and calibration against historical baseline years. Our per-word $\chi^2$/FDR layer is an addition, so this study should be read as a document-frequency variant with explicit testing rather than as a replication of their pipeline. We also apply Laplace smoothing symmetrically to both corpora, a natural extension of their single-corpus formula without a direct counterpart in the original.

\subsection{Re-Analysis of the Prior 651-Word List}
\label{sec:reanalysis}

The word universe is the 1{,}684 words tested in the prior study, of which 651 were reported significant under token-level $\chi^2$ with Bonferroni correction (458 in the AI-excess direction; 193 in the deficit direction, i.e., words humans use more). A prior word ``survives'' if it is significant under the new document-based criterion \emph{and} its direction ($r>1$ vs.\ $r<1$) matches the sign of the old excess score. This quantifies how much of the previously reported result was an artifact of token pseudoreplication.

\subsection{Rhythm Metrics}
\label{sec:rhythmmethod}

We compute 13 interpretable rhythm metrics per document. Markdown is preprocessed by removing code blocks, headings, tables, rules, images, and URLs, keeping the inner text of links and inline code; list items are kept as sentences with markers stripped. Sentences are split on 。！？ and newlines, following the \texttt{natural-japanese} lint~\citep{mizoguchi2026natural}; sentences without Japanese characters are dropped, and documents with fewer than 5 sentences are excluded.

\begin{itemize}
    \item \textbf{Sentence length (characters)}: mean, SD, and coefficient of variation (CV) of sentence lengths in characters.
    \item \textbf{Sentence length (moras)}: the same three statistics on approximated mora counts. Moras are approximated from the UniDic pronunciation field (katakana): small kana (ャュョァィゥェォヮ) merge into the preceding mora; ー, ッ, and ン each count as one mora; tokens without a reading (e.g., ASCII identifiers) contribute zero moras.
    \item \textbf{Burstiness} of the sentence-length sequence, $B = (\sigma - \mu)/(\sigma + \mu)$ with population SD~\citep{goh2008burstiness}, computed for character lengths and mora lengths. $B$ approaches $-1$ for regular (monotone) sequences and grows with bursty variation; the application to sentence-length sequences follows \texttt{natural-japanese}.
    \item \textbf{Paragraph sentence-count CV}: variation in the number of sentences per paragraph.
    \item \textbf{Noun-ending rate} (体言止め): fraction of prose sentences (list items excluded) whose final content morpheme is a noun.
    \item \textbf{ではなく per sentence}: rate of the contrastive construction ではなく (\textit{de wa naku}).
    \item \textbf{Commas per sentence}.
    \item \textbf{Sentence-ending bigram MATTR}: moving-average type--token ratio~\citep{covington2010mattr} over the sequence of sentence-final morpheme bigrams in prose sentences, window 20. Plain TTR depends on document length; MATTR removes this dependence. For documents with fewer than 20 ending bigrams (247 of 1{,}046) the window shrinks to the sequence length (recorded per document); fewer than 5 yields a missing value (15 documents).
\end{itemize}

Human/AI differences are summarized by two-sample Cohen's $d$ with pooled SD at the document level (positive $=$ AI higher), with bootstrap 95\% CIs ($n=2{,}000$); two-sided Mann--Whitney $U$ tests with Holm correction are reported as a secondary check, since rank tests are sensitive to distribution-shape differences at small $d$. These inferential statistics cover the 12 fixed-definition metrics; the ending-bigram MATTR, whose window shrinks for short documents, is reported descriptively (all-document and window-satisfying-subset estimates) without a CI or rank test.

\subsection{H1: Discrimination Design}
\label{sec:h1method}

We compare logistic-regression classifiers (L2, $C=1.0$) on three feature sets: (a)~lexical only (binary document occurrence over the full 25{,}902-word vocabulary), (b)~rhythm only (13 metrics), and (c)~both. Evaluation is stratified 5-fold cross-validation repeated 20 times ($100$ folds), with identical splits shared across feature sets so that AUC differences are paired at the fold level.

\paragraph{Leakage prevention.} Lexical feature selection runs \emph{inside each training fold} only: a minimum-document filter (present in $\geq 5$ and $\leq n-5$ training documents), the document-count $\chi^2$ test of Section~\ref{sec:lexmethod}, and BH-FDR $q<0.05$. The globally selected significant-word set is never used for classification, since selecting features on the full data would bias AUC upward. Selected vocabulary sizes are roughly 1{,}300--1{,}600 words per fold. Rhythm-feature imputation (training-fold median) and z-scoring likewise use training-fold statistics only.

\paragraph{Length control.} AI documents in this corpus are shorter than human articles, so any feature correlated with length could ride on a length confound. In the \texttt{lenadj} condition, both lexical and rhythm features are residualized against length covariates ($\log(1+\text{sentences})$, $\log(1+\text{characters})$) via linear regression fit within the training fold; a classifier using the two length covariates alone (\texttt{length\_only}) provides the reference.

\paragraph{Sensitivity analyses.} Llama 3.2 1B occasionally produces degenerate outputs (extreme sentence-length statistics). We report (i)~\texttt{sens\_degen}, excluding the 6 Llama documents exceeding the human 99th percentile on any of three length/comma statistics, and (ii)~\texttt{sens\_nollama}, excluding all 50 Llama documents.

\paragraph{Uncertainty.} 95\% CIs for AUCs and paired AUC differences are percentile bootstrap over the 100 fold values ($B=10{,}000$). Cross-validation folds share training data, so these CIs are approximate (anti-conservative); repeat-level means and SDs are reported alongside in the results files.

\subsection{H2: Fingerprint Layering Design}
\label{sec:h2method}

H2 asks whether rhythm fingerprints are homogeneous across models while lexical fingerprints differentiate them. We operationalize this at two levels, using the 350 AI documents.

\paragraph{Single-feature dispersion.} For each feature, the between/within ratio is the SD of the 7 per-model document means divided by the mean of the within-model SDs; smaller values mean more cross-model homogeneity. Features are the 13 rhythm metrics and the document-occurrence indicators of the 797 FDR-significant words. Because the ratio's null distribution differs between binary and continuous features, we run per-feature permutation tests (model labels shuffled, $n=1{,}000$, $+1$-corrected $p$-values~\citep{phipson2010permutation}) and also standardize each observed ratio against its permutation null ($z$). Lexical vs.\ rhythm ratios are compared with Mann--Whitney $U$ and rank-biserial correlation on both the raw ratio and $z$, with robustness variants (core-11 rhythm subset, words with $\geq 10$ AI documents, Llama excluded).

\paragraph{Multivariate model identification.} Single binary word indicators carry structurally large within-model variance (Bernoulli $p(1-p)$), which can understate the model-specificity of the lexical \emph{profile} as a whole. As a complementary multivariate probe, we train multinomial logistic regression to identify which of the 7 models produced a document (stratified 5-fold $\times$ 5 repeats; lexical features filtered only by within-fold frequency, no model labels used in selection), comparing lexical against rhythm feature sets; chance is $1/7 = 14.3\%$.

\section{Results}

\subsection{Document-Level Excess Vocabulary and the Correction of the Prior List}
\label{sec:lexresults}

Of the 1{,}684 tested words, 797 are significant under the document-based criterion (BH-FDR $q<0.05$); 381 additionally pass Bonferroni. Table~\ref{tab:survival} shows the fate of the prior study's 651 significant words.

\begin{table}[H]
\centering
\caption{Survival of the 651 words reported significant by the prior token-based analysis~\citep{imoto2026excess} under document-occurrence re-analysis. ``Survives'' requires significance under the new criterion with matching direction. Excess $=$ AI-overused; deficit $=$ human-overused.}
\label{tab:survival}
\small
\begin{tabular}{lrrr}
\toprule
Criterion & Survivors / 651 & Excess (/458) & Deficit (/193) \\
\midrule
Document $\chi^2$ + BH-FDR (primary) & \textbf{424 (65.1\%)} & 237 (51.7\%) & 187 (96.9\%) \\
Document $\chi^2$ + Bonferroni       & 259 (39.8\%) & 96 (21.0\%) & 163 (84.5\%) \\
Bootstrap 95\% CI excludes $r=1$ (uncorrected) & 458 (70.4\%) & 266 (58.1\%) & 192 (99.5\%) \\
All three criteria                   & 259 (39.8\%) & --- & --- \\
\bottomrule
\end{tabular}
\end{table}

Under the primary criterion, 35\% of the previously reported discoveries (227 of 651) are not supported at the document level; on the AI-excess side, roughly half vanish (458 $\to$ 237). Deficit-side words---high-frequency words such as 思う (\textit{omou}, ``to think''), which humans use far more often---survive almost entirely (96.9\%), because their differences are backed by broad document counts. Conversely, 354 words that were \emph{not} in the prior 651 are significant under the new criterion: document-level counting reveals words spread thinly across many documents that token counting under-weighted.

The failure mode of the token-based test is concrete. The prior study's top-ranked word ``Hash'' (excess score $+189.1$) occurred 402 times in the AI corpus, but in only 2 of 350 AI documents (vs.\ 2 of 700 human documents). Token-level $\chi^2$ turned within-document repetition into hundreds of pseudo-independent trials and returned $p \approx 0$; at the document level the word is indistinguishable from chance ($r=2.00$, 95\% CI $[0.50, 7.99]$, $q=0.955$). Of the prior top-5 words, four drop out under the document-based test in the same way; only レート (\textit{rate}; 16 vs.\ 11 documents, $r=2.83$, $q=0.018$) survives.

\begin{table}[H]
\centering
\caption{Top 10 excess words by frequency ratio $r$ and by frequency gap $\delta$ (FDR-significant, excess direction, document-occurrence statistics with Laplace smoothing). The bracketed 95\% CI accompanies the ranking index of each block: $r$ in the top block, $\delta$ in the bottom block. Doc counts are raw (AI /350, human /700); $r$ and $\delta$ use smoothed frequencies. Daggers ($\dagger$) mark ASCII fragments arising from MeCab segmentation of English spans embedded in Japanese text (see Limitations). Full statistics for all 1{,}684 words: \texttt{results/lexical/word\_stats.csv}.}
\label{tab:topwords}
\small
\begin{tabular}{llrrrr}
\toprule
 & Word & $r$ & $\delta$ & AI / Human docs & $q_{\text{FDR}}$ \\
\midrule
\multirow{10}{*}{by $r$}
 & NF$^\dagger$ & 21.97 [8.99, 59.91] & 0.060 & 21 / 1 & $1.6\times10^{-8}$ \\
 & スクリプティング (scripting) & 17.97 [7.99, 63.91] & 0.073 & 26 / 2 & $5.6\times10^{-10}$ \\
 & Found$^\dagger$ & 16.98 [6.49, 45.93] & 0.046 & 16 / 1 & $2.2\times10^{-6}$ \\
 & CSP & 16.98 [6.49, 47.93] & 0.046 & 16 / 1 & $2.2\times10^{-6}$ \\
 & クロスサイトリクエストフォージェリ & 13.31 [5.99, 49.93] & 0.053 & 19 / 2 & $5.3\times10^{-7}$ \\
 & Internal$^\dagger$ & 12.98 [4.66, 35.95] & 0.034 & 12 / 1 & $9.9\times10^{-5}$ \\
 & チェックリスト (checklist) & 12.38 [6.24, 37.95] & 0.081 & 30 / 4 & $2.2\times10^{-10}$ \\
 & デプロイメント (deployment) & 11.48 [5.19, 47.93] & 0.060 & 22 / 3 & $1.2\times10^{-7}$ \\
 & Foreign$^\dagger$ & 10.98 [3.99, 31.95] & 0.028 & 10 / 1 & $6.6\times10^{-4}$ \\
 & エスケープ (escape) & 10.78 [4.99, 30.98] & 0.070 & 26 / 4 & $1.1\times10^{-8}$ \\
\midrule
\multirow{10}{*}{by $\delta$}
 & 向上 (improvement) & 2.94 & 0.416 [0.359, 0.476] & 220 / 149 & $9.6\times10^{-38}$ \\
 & プラクティス (practice) & 6.51 & 0.393 [0.337, 0.447] & 162 / 49 & $1.2\times10^{-47}$ \\
 & ベスト (best) & 6.30 & 0.386 [0.330, 0.444] & 160 / 50 & $3.9\times10^{-46}$ \\
 & 纏め (summary) & 2.07 & 0.361 [0.307, 0.423] & 244 / 235 & $2.1\times10^{-26}$ \\
 & 実践 (put into practice) & 5.14 & 0.360 [0.308, 0.415] & 156 / 60 & $2.8\times10^{-39}$ \\
 & 品質 (quality) & 3.35 & 0.346 [0.288, 0.407] & 172 / 102 & $6.0\times10^{-31}$ \\
 & 効率 (efficiency) & 2.52 & 0.333 [0.275, 0.393] & 193 / 153 & $3.5\times10^{-25}$ \\
 & 解説 (explanation) & 2.37 & 0.306 [0.246, 0.364] & 185 / 156 & $1.9\times10^{-21}$ \\
 & 活用 (utilization) & 2.21 & 0.287 [0.225, 0.347] & 183 / 165 & $7.8\times10^{-19}$ \\
 & エンジニア (engineer) & 1.95 & 0.267 [0.205, 0.327] & 192 / 197 & $1.4\times10^{-15}$ \\
\bottomrule
\end{tabular}
\end{table}

Table~\ref{tab:topwords} lists the top excess words under each index. The two indices capture different strata, exactly as \citet{kobak2025excess} designed for English: $r$ surfaces low-frequency, AI-particular words (katakana technical loanwords such as スクリプティング, チェックリスト, デプロイメント), while $\delta$ surfaces high-frequency, function-adjacent vocabulary (向上, ベスト/プラクティス, 実践, 品質, 効率), the formal ``improvement talk'' register of instruction-tuned models. The $r$/$\delta$ complementarity thus transfers to Japanese.

\subsection{Rhythm: AI Text Is Monotone, in Every Model}
\label{sec:rhythmresults}

Table~\ref{tab:rhythm} reports the 13 rhythm metrics. The four largest human--AI differences are all \emph{variation-suppression} effects: burstiness of character-length sequences ($d=-0.96$), burstiness of mora-length sequences ($d=-0.80$), paragraph sentence-count CV ($d=-0.67$), and sentence-length CV ($d=-0.56$). AI text keeps sentence lengths and paragraph sizes in a narrow band where human writing swings; the mora-based replication of the character-based result indicates the monotony is phonological rather than an artifact of orthography.

\begin{table}[H]
\centering
\caption{Rhythm metrics, human (n=696) vs.\ AI (n=350), document level. $d$ = Cohen's $d$ (positive = AI higher) with bootstrap 95\% CI; MW $p$ = Mann--Whitney $U$ two-sided with Holm correction (secondary; rank tests are sensitive to shape differences at small $d$). The prior sentence-ending bigram TTR is replaced by MATTR (last row; see text).}
\label{tab:rhythm}
\small
\begin{tabular}{lrrrr}
\toprule
Metric & Human mean$\pm$SD & AI mean$\pm$SD & $d$ [95\% CI] & MW $p$ (Holm) \\
\midrule
Burstiness (chars) & $-0.223 \pm 0.167$ & $-0.382 \pm 0.162$ & $\mathbf{-0.96}$ [$-1.12, -0.81$] & $1.8\times10^{-51}$ \\
Burstiness (moras) & $-0.244 \pm 0.110$ & $-0.351 \pm 0.172$ & $\mathbf{-0.80}$ [$-1.00, -0.62$] & $2.1\times10^{-30}$ \\
Paragraph sent.-count CV & $0.662 \pm 0.273$ & $0.487 \pm 0.243$ & $\mathbf{-0.67}$ [$-0.83, -0.50$] & $8.2\times10^{-26}$ \\
Sentence-length CV (chars) & $0.679 \pm 0.344$ & $0.484 \pm 0.369$ & $-0.56$ [$-0.78, -0.36$] & $1.8\times10^{-51}$ \\
Sentence-length CV (moras) & $0.621 \pm 0.157$ & $0.521 \pm 0.380$ & $-0.39$ [$-0.83, -0.17$] & $2.1\times10^{-30}$ \\
Noun-ending rate (prose) & $0.137 \pm 0.183$ & $0.091 \pm 0.144$ & $-0.27$ [$-0.38, -0.15$] & $1.2\times10^{-16}$ \\
Sentence-length mean (chars) & $51.40 \pm 46.80$ & $40.06 \pm 75.93$ & $-0.19$ [$-0.48, -0.02$] & $2.8\times10^{-18}$ \\
Sentence-length mean (moras) & $40.81 \pm 23.07$ & $36.57 \pm 19.94$ & $-0.19$ [$-0.33, -0.07$] & $3.4\times10^{-6}$ \\
ではなく per sentence & $0.008 \pm 0.016$ & $0.011 \pm 0.021$ & $+0.17$ [$+0.04, +0.31$] & $0.40$ \\
Commas per sentence & $0.822 \pm 0.503$ & $1.104 \pm 2.894$ & $+0.16$ [$+0.08, +0.25$] & $0.075$ \\
Sentence-length SD (chars) & $44.06 \pm 84.77$ & $30.50 \pm 192.54$ & $-0.10$ [$-0.40, +0.05$] & $1.0\times10^{-78}$ \\
Sentence-length SD (moras) & $26.07 \pm 23.98$ & $23.30 \pm 75.41$ & $-0.06$ [$-0.40, +0.09$] & $3.5\times10^{-53}$ \\
Ending-bigram MATTR (w=20) & $0.703 \pm 0.141^{a}$ & $0.558 \pm 0.196^{a}$ & $-0.31$ (all docs) & --- \\
\bottomrule
\multicolumn{5}{l}{\footnotesize $^{a}$Means shown for the window-20-satisfying subset (human $n=613$, AI $n=171$; subset $d=-0.94$);}\\
\multicolumn{5}{l}{\footnotesize \phantom{$^{a}$}the headline $d=-0.31$ covers all documents with shrunken windows (see text).}
\end{tabular}
\end{table}

\paragraph{Direction is model-universal.} Burstiness (chars), paragraph sentence-count CV, and sentence-length CV (chars) are lower than human in \emph{all seven models}; per-model burstiness effect sizes range from $d=-0.59$ to $-1.72$. The degree varies: GPT-3.5 and GPT-4o are the most monotone (burstiness $d \approx -1.4$ to $-2.2$), while Claude Sonnet 4, Claude Opus 4, and GPT-OSS 20B come closest to human variation ($d \approx -0.3$ to $-0.6$). But no model deviates in the human direction. This is the pattern the practitioner claim predicts, and it holds across vendors, scales, and generations.

\paragraph{A length artifact, caught and corrected.} An earlier iteration of this analysis measured sentence-ending bigram diversity as plain TTR and found AI \emph{more} diverse ($d=+0.24$). TTR falls with sequence length, and human documents average 80.9 sentences against 38.5 for AI, so the comparison favored AI mechanically. Replacing TTR with MATTR (window 20) flips the sign: $d=-0.31$ over all documents, and $d=-0.94$ within the window-20-satisfying subset (which over-represents long AI documents, so the subset value is subset-specific). AI sentence endings are less diverse, not more. We report this reversal as a caution for stylometric work: length-uncontrolled diversity metrics can invert conclusions.

\subsection{H1: Lexical Features Dominate Discrimination}
\label{sec:h1results}

Table~\ref{tab:auc} reports human/AI classification AUCs. Lexical features reach AUC 0.998; the 13 rhythm metrics alone reach 0.897. The paired difference is $-0.101$ (95\% CI $[-0.105, -0.097]$); the lexical classifier wins in 100 of 100 folds. H1 is supported---discriminative power differs by feature family---but in the direction \emph{opposite} to the practitioner claim: vocabulary, not rhythm, carries most of the human/AI signal in this corpus.

\begin{table}[H]
\centering
\caption{Human/AI discrimination AUC (mean $\pm$ SD over 100 folds; brackets: bootstrap 95\% CI, $B=10{,}000$). \texttt{main}: all 1{,}046 documents. \texttt{lenadj}: all features residualized against document-length covariates. \texttt{sens\_degen}: 6 degenerate Llama documents removed. \texttt{sens\_nollama}: all 50 Llama documents removed.}
\label{tab:auc}
\small
\begin{tabular}{lccc}
\toprule
Condition & Lexical & Rhythm & Combined \\
\midrule
main & $0.998 \pm 0.003$ [0.997, 0.998] & $0.897 \pm 0.020$ [0.893, 0.901] & $0.998 \pm 0.003$ [0.998, 0.999] \\
lenadj & $0.974 \pm 0.011$ [0.972, 0.976] & $0.810 \pm 0.031$ [0.804, 0.816] & $0.980 \pm 0.009$ [0.978, 0.981] \\
sens\_degen & $0.998 \pm 0.003$ [0.997, 0.998] & $0.913 \pm 0.020$ [0.909, 0.917] & $0.998 \pm 0.003$ [0.998, 0.999] \\
sens\_nollama & $0.999 \pm 0.003$ [0.998, 0.999] & $0.921 \pm 0.018$ [0.918, 0.924] & $0.999 \pm 0.002$ [0.998, 0.999] \\
\midrule
\multicolumn{4}{l}{\texttt{lenadj} reference, length covariates only: $0.811 \pm 0.024$ [0.806, 0.815]} \\
\bottomrule
\end{tabular}
\end{table}

\begin{table}[H]
\centering
\caption{Paired AUC differences (same folds), with bootstrap 95\% CIs and the fraction of the 100 folds in which the difference is positive.}
\label{tab:aucdiff}
\small
\begin{tabular}{llrr}
\toprule
Condition & Comparison & Difference [95\% CI] & Folds $>0$ \\
\midrule
main & rhythm $-$ lexical & $\mathbf{-0.101}$ [$-0.105, -0.097$] & 0/100 \\
main & combined $-$ lexical & $+0.0003$ [$+0.0002, +0.0004$] & 67/100 \\
main & combined $-$ rhythm & $+0.101$ [$+0.097, +0.105$] & 100/100 \\
lenadj & rhythm $-$ lexical & $-0.164$ [$-0.170, -0.157$] & 0/100 \\
lenadj & combined $-$ lexical & $+0.006$ [$+0.005, +0.007$] & 89/100 \\
lenadj & lexical $-$ length\_only & $+0.163$ [$+0.159, +0.168$] & 100/100 \\
lenadj & rhythm $-$ length\_only & $-0.000$ [$-0.008, +0.008$] & 46/100 \\
sens\_degen & rhythm $-$ lexical & $-0.085$ [$-0.088, -0.081$] & 0/100 \\
sens\_nollama & rhythm $-$ lexical & $-0.078$ [$-0.081, -0.074$] & 0/100 \\
\bottomrule
\end{tabular}
\end{table}

Four readings of Tables~\ref{tab:auc} and \ref{tab:aucdiff}:

\begin{enumerate}
    \item \textbf{The lexical advantage is robust.} It survives length residualization ($-0.164$), degenerate-document removal ($-0.085$), and full Llama exclusion ($-0.078$); no CI approaches zero.
    \item \textbf{Rhythm alone is still practically useful.} Thirteen interpretable numbers reach AUC 0.897 (0.913--0.921 after outlier handling), close to nine-tenths discrimination with features a human can read.
    \item \textbf{Length explains much of the rhythm signal.} Document length alone discriminates at AUC 0.811 (AI documents are shorter). After residualization, rhythm reaches 0.810, matching the length-only baseline rather than exceeding it (paired difference $-0.000$ [$-0.008, +0.008$]). Since the residualized features are (linearly) orthogonal to length, this 0.810 is carried by length-independent variation; but rhythm does not outperform knowing length alone in this setup, and its headline AUC of 0.897 should not be read as length-free. Lexical features are nearly untouched by residualization (0.974).
    \item \textbf{Rhythm is non-redundant on top of lexicon.} In \texttt{main} the combined gain is $+0.0003$: real but immaterial, since lexical AUC sits at a 0.998 ceiling. Under length adjustment the gain is $+0.006$ [$+0.005, +0.007$], positive in 89/100 folds: rhythm contributes information the vocabulary does not contain.
\end{enumerate}

\subsection{H2: A Two-Layer Fingerprint}
\label{sec:h2results}

\paragraph{Single-feature level: rhythm degrees differ across models.}
All 13 rhythm metrics differentiate the 7 models (permutation $p<0.005$ for every metric), with between/within ratios from 0.32 to 1.23 (Table~\ref{tab:ratios}). Single significant words are, on average, \emph{less} model-differentiating: median ratio 0.27 (IQR 0.20--0.37), median permutation $z=2.1$, with only 55.2\% of the 797 words reaching permutation $p<0.05$. The lexical-vs-rhythm contrast is significant and large in the direction opposite to naive expectation (Mann--Whitney $p=2.0\times10^{-7}$, rank-biserial $r=-0.84$ on ratios; $p=5.0\times10^{-7}$, $r=-0.81$ on permutation $z$), and is stable under the core-11 rhythm subset, restriction to words in $\geq 10$ AI documents, and Llama exclusion ($r=-0.73$ to $-0.86$). Part of this is structural: binary occurrence indicators carry Bernoulli within-model variance that depresses per-word ratios.

\begin{table}[H]
\centering
\caption{Between/within model dispersion of rhythm metrics across the 7 LLMs (AI documents only). All permutation $p<0.005$ ($n=1{,}000$, $+1$-corrected). The two sentence-ending rhetoric metrics (bottom) separate cleanly from the core rhythm metrics.}
\label{tab:ratios}
\small
\begin{tabular}{lrr}
\toprule
Metric & Ratio & Perm.\ $z$ \\
\midrule
ではなく per sentence & 0.32 & 4.3 \\
Sentence-length CV (chars) & 0.44 & 5.2 \\
Sentence-length CV (moras) & 0.47 & 6.0 \\
Burstiness (chars) & 0.52 & 9.4 \\
Sentence-length mean (chars) & 0.53 & 8.3 \\
Sentence-length SD (chars) & 0.57 & 4.9 \\
Paragraph sent.-count CV & 0.58 & 10.6 \\
Burstiness (moras) & 0.58 & 10.6 \\
Commas per sentence & 0.62 & 7.7 \\
Sentence-length SD (moras) & 0.64 & 5.6 \\
Sentence-length mean (moras) & 0.65 & 10.4 \\
\midrule
\textbf{Noun-ending rate (prose)} & \textbf{0.85} & \textbf{16.9} \\
\textbf{Ending-bigram MATTR} & \textbf{1.23} & \textbf{26.7} \\
\midrule
\multicolumn{3}{l}{Lexical (797 FDR words): median ratio 0.27 (IQR 0.20--0.37), median $z=2.1$}\\
\bottomrule
\end{tabular}
\end{table}

\paragraph{Multivariate level: the lexical profile names the model.}
The picture reverses when features act together. A multinomial classifier over lexical document-occurrence features identifies which of the 7 models produced a document with $92.1\% \pm 3.8\%$ accuracy (macro-F1 0.920; chance 14.3\%); the 13 rhythm metrics manage $50.9\% \pm 6.5\%$ (macro-F1 0.495). Excluding Llama (6-way, chance 16.7\%): lexical 90.8\%, rhythm 50.2\%. Individually weak word indicators are collectively a near-unique model signature; rhythm metrics, individually well-separated in degree, do not accumulate into one.

\paragraph{Model-specific vocabulary is concrete and interpretable.} The most model-differentiating words are ordinary vocabulary choices rather than exotica: 纏め (\textit{matome}, ``summary''; ratio 1.59) appears in essentially every Claude Opus 4 and GPT-OSS document but 2\% of Llama documents (GPT-3.5: 34\%); 於く (\textit{oku}, formal ``in/at''; 1.25) ranges from 2\% (GPT-OSS) to 98\% (Sonnet 4); 用いる (\textit{mochiiru}, ``to employ''; 1.24) is nearly GPT-4o-exclusive (42\% vs.\ 0--6\% elsewhere).

\paragraph{Sentence-ending rhetoric behaves like vocabulary.} Within the rhythm family, the two sentence-ending metrics are outliers: noun-ending rate (ratio 0.85) and ending-bigram MATTR (1.23) top the dispersion table, well separated from the core rhythm band (0.32--0.65). Claude 3 Haiku, GPT-3.5, and GPT-4o use noun endings almost never (0.3--1.3\%), while Claude Sonnet 4, Opus 4, and GPT-OSS 20B use them at or above the human rate (13--22\% vs.\ human 13.7\%). With only two such metrics we treat this as a descriptive observation rather than a tested contrast, but it suggests sentence-final rhetoric belongs to the model-specific layer.

\paragraph{H2 verdict.} The hypothesis splits by level of operationalization. At the single-feature level it fails: every rhythm metric differs across models in \emph{degree} (all permutation $p<0.005$). At the multivariate level it holds: lexical profiles identify the model (92\%), rhythm profiles largely do not (51\%). The defensible summary is that the rhythm shift is \emph{directionally} homogeneous---all 7 models are more monotone than humans on the core metrics---while its magnitude is model-dependent; the lexical profile is the model-specific fingerprint.

\section{Discussion}

\subsection{A Two-Layer Fingerprint Structure}

The results organize into two layers. \textbf{Layer 1 (shared machine accent)}: rhythmic monotonization (suppressed burstiness, uniform paragraph structure, compressed sentence-length variation) points the same way in all 7 models and supports human/AI discrimination at AUC $\approx 0.90$ with 13 interpretable features. \textbf{Layer 2 (model signature)}: the lexical profile, which identifies the generating model at 92\% accuracy, plus sentence-ending rhetoric, which disperses across models more than any core rhythm metric. Human/AI discrimination is best served by vocabulary (AUC 0.998), which aggregates both layers; model attribution is served almost solely by layer 2. Compactly: \emph{what} is written tells you which model; \emph{how} it is written tells you a machine wrote it.

This structure connects to the homogenization literature: \citet{munozortiz2024contrasting} observed for English that inter-LLM differences are smaller than the human--LLM gap on dispersion measures, and our layer-1 finding replicates that shape in Japanese at effect sizes up to $d=-0.96$. That the shared layer is rhythmic and the differentiating layer is lexical is also consistent with alignment-tuning accounts~\citep{juzek2025overuse, rlhf2025detectability, padmakumar2024diversity}: pressure toward well-organized, uniformly readable output plausibly acts on all tuned models alike, while each model's lexical preferences remain idiosyncratic to its training mix and feedback data.

\subsection{The Practitioner Claim, Tested}

The \texttt{natural-japanese} lint~\citep{mizoguchi2026natural} deserves credit for the sharpest form of the rhythm intuition, and our data vindicate its descriptive core: the largest human--AI rhythm gaps are exactly the quantities it lints (burstiness, paragraph-structure CV), the direction is universal across models, and the mora-level replication shows the effect is phonological, not orthographic. Where the data part ways with the claim is the comparative statement. On this corpus, vocabulary out-discriminates rhythm by a wide, stable margin ($\Delta$AUC $\approx 0.10$, never reversed in 100 folds, robust to length control and model exclusions). ``AI smell lives in rhythm rather than vocabulary'' should be revised to: rhythm is where AI text \emph{feels} uniform across models, but vocabulary is where it is most \emph{detectable}, at least in the technical-blog register tested here. The practical case for rhythm features is different and still strong: they are few, cheap, human-readable, and directionally universal, which suits lint-style feedback tools precisely because writers can act on ``vary your sentence lengths'' but not on ``use 向上 less across 800 words.''

\subsection{Correcting Our Own Baseline}

This study's lexical layer began as an audit of our prior work, and the audit found a real defect: token-level $\chi^2$ tests treat within-document repetition as independent evidence, and 35\% of the 651 words we previously reported---about half of the AI-excess side---do not survive document-based re-analysis. We prefer stating this plainly over defending the earlier number: the corrected list (237 excess words surviving; 797 significant words in the full re-analysis) is smaller and better-founded, and the failure cases (``Hash'': 402 tokens, 2 documents) illustrate why document occurrence is the right unit for excess-vocabulary claims in any language. The correction also sharpens the cross-linguistic result: the $r$/$\delta$ complementarity of \citet{kobak2025excess} (ratio for rare AI-particular words, gap for high-frequency register words) emerges in Japanese only once the statistics are computed at the document level.

\section{Limitations}
\label{sec:limitations}

\paragraph{Register and topic confound.} The human corpus consists of real, freely composed Qiita/Zenn articles; the AI corpus was generated from fixed prompts on 10 themes. Topic mix, prompt framing, and register therefore differ between classes, and the lexical AUC of 0.998 reflects within-corpus separation rather than operational AI-detector performance. The words driving separation are ordinary (思う: 75\% of human vs.\ 6\% of AI documents; 向上: 21\% vs.\ 63\%), and we verified that no single artifact token fully separates the classes, but prompt style and topic contributions cannot be excised. All H1 claims are scoped to this same-theme-family corpus contrast.

\paragraph{Approximate fold CIs.} Cross-validation folds share training data, so bootstrap CIs over fold AUCs are anti-conservative approximations. Repeat-level aggregates are provided in the released results; the headline contrast (0/100 folds reversed, present in all four conditions) does not depend on the CI construction.

\paragraph{Mora approximation.} Mora counts derive from UniDic pronunciation fields; tokens without readings (ASCII identifiers, code fragments) contribute zero moras, deflating mora-based lengths for English-heavy documents. Character-based metrics are reported alongside throughout, and the burstiness findings agree across the two bases.

\paragraph{Tokenization artifacts.} Several high-$r$ items are ASCII fragments (NF, Found, Internal, Foreign) produced by MeCab segmentation of embedded English spans, a limitation inherited from the shared tokenizer of both corpora; we flag them in Table~\ref{tab:topwords} rather than silently exclude them.

\paragraph{Inherited corpus constraints.} The source corpus was generated at default temperature without fixed seeds, so corpus-level regeneration is not bit-reproducible (all analyses on the fixed corpus are seeded and deterministic). MATTR windows shrink for short documents (247 of 1{,}046, disproportionately AI), so window widths are not uniform; both the all-document and window-satisfying estimates are reported. The H2 lexical set was selected on a human/AI criterion, which is not leakage for the model-dispersion question but is not guaranteed orthogonal to it; restriction to broadly attested words leaves the conclusion unchanged. The combined classifier concatenates binary and z-scored features under a single L2 penalty without per-block tuning, and length residualization is linear, leaving possible nonlinear length effects. The sentence-ending rhetoric group contains only two metrics and is reported descriptively.

\paragraph{Scope.} One language (Japanese), one domain (technical blogging), seven models. Our prior work found excess vocabulary to be strongly domain-dependent~\citep{imoto2026excess}, so the lexical/rhythm balance may differ in other genres; the two-layer structure is the hypothesis we would test there.

\section{Conclusion}

We compared lexical and rhythmic fingerprints of Japanese LLM-generated text across 7 models under a controlled design, after rebuilding the lexical statistics on a document-occurrence basis. Three results stand. First, the document-based re-analysis corrects our prior study: 35\% of its 651 reported excess words---half of the AI-excess side---were artifacts of token-level pseudoreplication, and the corrected Japanese excess vocabulary reproduces the $r$/$\delta$ complementarity established for English. Second, rhythmic monotonization is a genuine, model-universal property of Japanese AI text (burstiness $d=-0.96$; direction shared by all 7 models), yet lexical features dominate human/AI discrimination (AUC 0.998 vs.\ 0.897; gap robust to length control and outlier exclusion), the opposite of the practitioner claim that motivated the comparison, whose descriptive core our rhythm measurements still confirm. Third, the two families form a two-layer fingerprint: rhythm shifts in a shared direction with model-dependent magnitude, while the lexical profile (and, descriptively, sentence-ending rhetoric) identifies the generating model at 92\% accuracy. Vocabulary tells you which model wrote a text; rhythm tells you that a model wrote it.

\paragraph{Data Availability.}
All analysis code, per-word statistics (1{,}684 words), rhythm measurements (1{,}046 documents), and classification results are available at \url{https://github.com/kenimo49/llm-text-fingerprints-ja} under the MIT License.
% TODO: Zenodo concept DOI to be inserted after deposit (placeholder: 10.5281/zenodo.XXXXXXX)
The underlying corpora are those of \citet{imoto2026excess} (\url{https://github.com/kenimo49/excess-vocabulary-ja}), used read-only at commit \texttt{fdb24f4c}.

\bibliography{references}

\end{CJK}
\end{document}
